\section{Dinâmica acoplada}
\label{sec_dinamica_acoplada_rev_bibliografica}


Como apresentado por \cite{adachi1999}, experimentos de identificação de sistema em órbita foram conduzidos utilizando o satélite japonês ETS-VI (Engineering Test Satellite VI).
O objetivo foi verificar a consistência de modelos dinâmicos previamente obtidos por testes em solo, por meio da análise de dados de telemetria de atitude e acelerômetros durante excitações controladas aplicadas à espaçonave.

Posteriormente, \cite{yoshida2003} apresentou a missão ETS-VII (Engineering Test Satellite VII), que deu sequência a demonstrações em órbita ao incorporar um manipulador robótico, permitindo investigar o acoplamento dinâmico entre a atitude da espaçonave e a reação do manipulador durante manobras robóticas.

Wei et al.~\cite{wei2011} apresentaram uma lei de controle acoplada para coordenar posição relativa e atitude no acoplamento autônomo com um alvo em rotação livre, assegurando o alinhamento das portas de acoplamento ao longo da aproximação em \gls{rvdb}.